%% toc: P15 159 159.7,159.10
\chapter{Regular Conditional Distributions Defined Almost Everywhere}
\label{sec:p15-w02b-p15-008-ae-kernels}

\Cref{def:p15-regular-conditional-law} asks a regular conditional law to
be a probability measure at every conditioning argument, and
\cref{thm:p15-kernel-expectation} integrates functions against such a
kernel. Durrett's definition of a regular conditional distribution is
weaker: the measure condition is required only for almost every outcome
\cite[p.~214]{DurrettPTEAuthor2019}. This chapter proves Durrett's
integration formula, his Theorem~4.1.16, in that weaker regime. It also
states exactly when such an object can be replaced by a kernel of
\cref{def:p15-regular-conditional-law}, and gives a model in which the
replacement by a change on a conditioning-null set is impossible.

\emph{Inputs.} A probability space, a measurable map into a measurable
space, a sub-sigma-algebra, and a function \(\mu\) satisfying the source's
two conditions. \emph{Conclusions.} For every measurable \(f\geq0\), and for
every real \(f\) with \(E|f(X)|<\infty\), integration of \(f\) against
\(\mu(\omega,\cdot)\) gives a conditional expectation of \(f(X)\) outside a
null set, and the exceptional set is controlled explicitly.

\section{The source definition}
\label{sec:p15-w02b-p15-008-source-definition}

The notation follows Durrett. The probability space is written
\((\Omega,\mathcal F,P)\); \(\mathcal F\) plays the role of
\(\mathcal F_0\) in \cref{def:p15-general-conditional-expectation}, and
\(C_{\mathcal G}\) is the version notation introduced there.

\begin{hbdefinition}{def:p15-w02b-p15-008-rcd}{A regular conditional distribution in the source's sense}
Let \((\Omega,\mathcal F,P)\) be a probability space, let
\(X:(\Omega,\mathcal F)\to(S,\mathcal S)\) be measurable, and let
\(\mathcal G\subseteq\mathcal F\) be a sigma-algebra.
A function \(\mu:\Omega\times\mathcal S\to[0,1]\) is a regular conditional
distribution for \(X\) given \(\mathcal G\) if
\begin{enumerate}
\item for each \(A\in\mathcal S\), the function
      \(\omega\mapsto\mu(\omega,A)\) is a version of
      \(P(X\in A\mid\mathcal G)\); in particular it is
      \(\mathcal G\)-measurable;
\item for almost every \(\omega\), the set function
      \(A\mapsto\mu(\omega,A)\) is a probability measure on
      \((S,\mathcal S)\).
\end{enumerate}
When \(S=\Omega\) and \(X\) is the identity map, \(\mu\) is called a
regular conditional probability \cite[p.~214]{DurrettPTEAuthor2019}.
\index{regular conditional distribution!almost-everywhere definition}
\end{hbdefinition}

\begin{hbconvention}{def:p15-w02b-p15-008-exceptional-set}{Reading of ``almost every'' and the exceptional set}
Condition~(ii) is read as follows: there is a set \(N\in\mathcal F\) with
\(P(N)=0\) such that \(\mu(\omega,\cdot)\) is a probability measure on
\((S,\mathcal S)\) for every \(\omega\in\Omega\setminus N\).
Such an \(N\) is called an exceptional set for \(\mu\).
The set of arguments at which \(\mu(\omega,\cdot)\) fails to be a
probability measure is a subset of \(N\). Nothing in the definition makes
\(N\), or that subset, a member of \(\mathcal G\); see
\cref{thm:p15-w02b-p15-008-ae-counterexample}.
\end{hbconvention}

Every regular conditional law \(K\) of
\cref{def:p15-regular-conditional-law} is a regular conditional
distribution in this sense with exceptional set \(N=\emptyset\): its
event functions are \(\mathcal G\)-measurable, agree almost everywhere with
versions of the event conditional probabilities and therefore are versions
themselves, and each \(K(\omega,\cdot)\) is a probability measure.
The source states its integration theorem with the conditioning
sigma-algebra written \(\mathcal F\); in the definition it is
\(\mathcal G\). Below it is \(\mathcal G\) throughout.

\section{Integration against an almost-everywhere kernel}
\label{sec:p15-w02b-p15-008-integration}

\begin{hbtheorem}{thm:p15-w02b-p15-008-ae-integration}{proved=annex:p15-w02b-p15-008-ae-integration}{Integrating a function against a source-sense conditional law}
Let \(\mu\) be a regular conditional distribution for \(X\) given
\(\mathcal G\), with exceptional set \(N\) as in
\cref{def:p15-w02b-p15-008-exceptional-set}.
\begin{enumerate}
\item For every measurable \(f:S\to[0,\infty]\) there is a
      \(\mathcal G\)-measurable \(g_f:\Omega\to[0,\infty]\) such that
      \[
           g_f(\omega)=\int\mu(\omega,dx)f(x)
                         \qquad(\omega\in\Omega\setminus N),
      \]
      and \(g_f\) is a version of the nonnegative conditional expectation
      \(C_{\mathcal G}f(X)\) of \cref{thm:p15-conditional-limits}.
      The set \(\Omega\setminus N\) on which the identity holds is the same
      for every \(f\).
\item Let \(f:(S,\mathcal S)\to(\mathbb R,\mathcal R)\) be measurable with
      \(E|f(X)|<\infty\). There is a set \(M_f\in\mathcal G\) with
      \(P(M_f)=0\) such that for every
      \(\omega\notin N\cup M_f\) the integral
      \(\int\mu(\omega,dx)|f(x)|\) is finite, and there is a version
      \(Z_f\) of \(E(f(X)\mid\mathcal G)\) with
      \[
           Z_f(\omega)=\int\mu(\omega,dx)f(x)
                         \qquad(\omega\notin N\cup M_f).
      \]
      Hence
      \[
           E(f(X)\mid\mathcal G)=\int\mu(\omega,dx)f(x)
                                 \quad\text{a.s.}
      \]
\end{enumerate}
\end{hbtheorem}

The displayed almost-sure formula is Durrett's Theorem~4.1.16, whose proof
runs from indicators through simple and nonnegative functions to
\(f=f^+-f^-\) \cite[p.~214]{DurrettPTEAuthor2019}. The annex follows that
route. The right side of the formula is defined only outside
\(N\cup M_f\), a set that need not belong to \(\mathcal G\). The theorem
therefore asserts two things: the right side coincides there with a
\(\mathcal G\)-measurable version, and every version agrees with it almost
surely. When \(N=\emptyset\), part~(i) is the nonnegative case of
\cref{thm:p15-kernel-expectation}.

\section{Conditioning-measurable exceptional sets}
\label{sec:p15-w02b-p15-008-reduction}

\begin{hbcorollary}{thm:p15-w02b-p15-008-ae-reduction}{proved=annex:p15-w02b-p15-008-ae-reduction}{Reduction to a kernel defined at every argument}
Let \(\mu\) be a regular conditional distribution for \(X\) given
\(\mathcal G\) with exceptional set \(N\). Suppose \(N\subseteq M_0\) for
some \(M_0\in\mathcal G\) with \(P(M_0)=0\). Fix \(s_0\in S\) and put
\[
 K(\omega,A)=
 \begin{cases}
  \mu(\omega,A),&\omega\notin M_0,\\
  \mathbf1_A(s_0),&\omega\in M_0.
 \end{cases}
\]
Then \(K\) is a regular conditional law of \(X\) given \(\mathcal G\) in the
sense of \cref{def:p15-regular-conditional-law}, and
\(K(\omega,\cdot)=\mu(\omega,\cdot)\) for every \(\omega\notin M_0\).
For measurable \(f\geq0\), and for real \(f\) with \(E|f(X)|<\infty\),
the versions supplied by \cref{thm:p15-kernel-expectation} for \(K\) and by
\cref{thm:p15-w02b-p15-008-ae-integration} for \(\mu\) agree almost surely,
and the two kernel integrals coincide at every argument outside
\(M_0\) at which they are defined.
The hypothesis holds with \(M_0=N\) whenever \(\mathcal G\) contains every
\(P\)-null set of \(\mathcal F\).
\end{hbcorollary}

The point \(s_0\) exists because \(\Omega\) is nonempty, as \(P(\Omega)=1\),
and \(X\) maps it into \(S\). The existence theorems
\cref{thm:p15-real-conditional-law,thm:p15-standard-conditional-law}
already produce kernels defined at every argument, with their common
exceptional sets chosen in \(\mathcal G\). They cover the source's
Theorems~4.1.17 and 4.1.18 \cite[p.~215]{DurrettPTEAuthor2019}.
\Cref{thm:p15-w02b-p15-008-ae-integration} is needed for a source-sense
object supplied in some other way, and it does not depend on the
hypothesis of the corollary.

\section{An exceptional set outside the conditioning sigma-algebra}
\label{sec:p15-w02b-p15-008-counterexample}

\begin{hbtheorem}{thm:p15-w02b-p15-008-ae-counterexample}{proved=annex:p15-w02b-p15-008-ae-counterexample}{A source-sense conditional law that no conditioning-null change repairs}
Let \(\Omega=S=[0,2]\), let \(\mathcal F=\mathcal S\) be the Borel
sigma-algebra of \([0,2]\), let \(\lambda\) be Lebesgue measure, and put
\(P(A)=\nu(A)=\lambda(A\cap[0,1])\) for \(A\in\mathcal F\). Let
\(\mathcal G\) be the family of subsets of \([0,2]\) that are countable or
have countable complement, and let \(X\) be the identity map. Define
\[
 \mu(\omega,A)=
 \begin{cases}
  1/2,&\omega\in(1,2]\text{ and }A=\{\omega\},\\
  \nu(A),&\text{otherwise}.
 \end{cases}
\]
Then:
\begin{enumerate}
\item \(\mathcal G\) is a sigma-algebra contained in \(\mathcal F\), and
      \(\mu\) is a regular conditional distribution for \(X\) given
      \(\mathcal G\) in the sense of \cref{def:p15-w02b-p15-008-rcd}.
      The set of \(\omega\) at which \(\mu(\omega,\cdot)\) is not a
      probability measure is exactly \(N=(1,2]\), and \(P(N)=0\).
\item No \(M\in\mathcal G\) with \(P(M)=0\) contains \(N\). Consequently
      there is no probability kernel \(K\) from \((\Omega,\mathcal G)\) to
      \((S,\mathcal S)\) and \(\mathcal G\)-measurable null set \(M\) with
      \(K(\omega,\cdot)=\mu(\omega,\cdot)\) for all \(\omega\notin M\).
\item The constant kernel \(K_0(\omega,A)=\nu(A)\) is a regular conditional
      law of \(X\) given \(\mathcal G\), and \(K_0(\cdot,A)=\mu(\cdot,A)\)
      outside a set of at most one point for each \(A\in\mathcal S\).
\item Let \(f\) be Borel on \([0,2]\) with \(E|f(X)|<\infty\), and put
      \(c_f=\int f\,d\nu\). The constant \(c_f\) is a version of
      \(E(f(X)\mid\mathcal G)\), and
      \(\int\mu(\omega,dx)f(x)=c_f\) for every \(\omega\in[0,1]\).
      The function equal to \(c_f\) on \([0,1]\) and to \(0\) on
      \((1,2]\) is \(\mathcal G\)-measurable if and only if \(c_f=0\).
\end{enumerate}
\end{hbtheorem}

\begin{hbexample}{def:p15-w02b-p15-008-worked}{Computing conditional moments from the source-sense law}
In the model of \cref{thm:p15-w02b-p15-008-ae-counterexample}, take
\(f(x)=x\) and \(f(x)=x^2\). Both satisfy \(E|f(X)|<\infty\), since \(X\)
takes values in \([0,1]\) almost surely. For every \(\omega\in[0,1]\),
\(\mu(\omega,\cdot)=\nu\) is the uniform law on \([0,1]\), so
\[
 \int\mu(\omega,dx)\,x=\int_0^1x\,dx=\frac12,\qquad
 \int\mu(\omega,dx)\,x^2=\int_0^1x^2\,dx=\frac13 .
\]
At \(\omega\in(1,2]\) the set function \(\mu(\omega,\cdot)\) is not a
measure and no integral is formed. Part~(ii) of
\cref{thm:p15-w02b-p15-008-ae-integration} gives
\(E(X\mid\mathcal G)=1/2\) and \(E(X^2\mid\mathcal G)=1/3\) almost surely,
in agreement with part~(iv) of the counterexample theorem. By
\cref{def:p15-conditional-variance},
\[
 \operatorname{Var}(X\mid\mathcal G)=\frac13-\frac14=\frac1{12}
                                      \quad\text{a.s.}
\]
The versions to report are the constants \(1/2\), \(1/3\) and \(1/12\).
The functions \(\frac12\mathbf1_{[0,1]}\) and \(\frac13\mathbf1_{[0,1]}\),
obtained by assigning \(0\) where the integral is undefined, equal these
almost surely but are not \(\mathcal G\)-measurable, so they are not
versions.
\end{hbexample}

\begin{hbwarning}{def:p15-w02b-p15-008-limits}{What the almost-everywhere definition does not give}
\begin{enumerate}
\item The kernel integral is defined only off the exceptional set. Filling
      it in by an arbitrary value there can destroy
      \(\mathcal G\)-measurability, as part~(iv) of
      \cref{thm:p15-w02b-p15-008-ae-counterexample} shows; this is the
      general fact, recorded in \cref{thm:p15-conditional-limits}, that
      modifications on an \(\mathcal F\)-null set need not preserve
      \(\mathcal G\)-measurability, which
      \cref{thm:p15-w02b-p15-008-null-modification} characterizes and
      witnesses with trivial \(\mathcal G\). The correct statement is the one in
      \cref{thm:p15-w02b-p15-008-ae-integration}: agreement with a
      \(\mathcal G\)-measurable version off a null set.
\item The extra null set \(M_f\) for a signed integrand depends on \(f\).
      Only the nonnegative identity of part~(i) holds on one set
      \(\Omega\setminus N\) for all integrands simultaneously.
\item Neither definition guarantees existence. The source records that
      regular conditional distributions do not always exist
      \cite[p.~214]{DurrettPTEAuthor2019}; existence is proved here only
      under the hypotheses of
      \cref{thm:p15-real-conditional-law,thm:p15-standard-conditional-law}.
\item Part~(iii) of the counterexample shows that a kernel of
      \cref{def:p15-regular-conditional-law} may still exist when the
      corollary's hypothesis fails. It then differs from \(\mu(\omega,\cdot)\)
      as a measure at uncountably many arguments, although for each fixed
      event the two event functions agree almost surely.
\end{enumerate}
\end{hbwarning}

\section{Null modifications and decreasing conditional limits}
\label{sec:p15-w02b-p15-008-version-limits}

The repairs above turn on which changes on a null set keep a function a
conditional version. This section answers that question for scalar
conditional expectations and records the decreasing form of conditional
monotone convergence. It uses the notation of
\cref{def:p15-general-conditional-expectation}: a probability space
\((\Omega,\mathcal F_0,P)\), a sub-sigma-algebra
\(\mathcal G\subseteq\mathcal F_0\), and \(C_{\mathcal G}\) for a selected
version. The letters \(Y_n\) and \(Y\) are real random variables, as in
the source's remark; they are not the random element \(X\) of the
preceding sections.

\emph{Inputs.} An integrable variable with one version of its conditional
expectation, or a sequence of variables decreasing almost everywhere.
\emph{Conclusions.} An exact criterion for a null-set modification to
remain a version, a model in which a modification on an
\(\mathcal F_0\)-null set fails it, and almost-everywhere convergence of
conditional expectations along decreasing sequences, with a model showing
that integrability of the first term cannot be dropped.

\begin{hbtheorem}{thm:p15-w02b-p15-008-null-modification}{proved=annex:p15-w02b-p15-008-null-modification}{Which null modifications keep a conditional version}
Let \(X\) be integrable and let \(Z\) be a version of
\(E(X\mid\mathcal G)\).
\begin{enumerate}
\item If \(D\in\mathcal G\), \(P(D)=0\), and \(W:\Omega\to\mathbb R\) is
      \(\mathcal G\)-measurable, then
      \(Z\mathbf1_{\Omega\setminus D}+W\mathbf1_D\) is a version of
      \(E(X\mid\mathcal G)\).
\item Let \(Z':\Omega\to\mathbb R\) satisfy \(Z'=Z\) outside some
      \(D'\in\mathcal F_0\) with \(P(D')=0\). Then \(Z'\) is a version of
      \(E(X\mid\mathcal G)\) if and only if \(Z'\) is
      \(\mathcal G\)-measurable.
\item Let \(\Omega=[0,1]\), let \(\mathcal F_0\) be its Borel
      sigma-algebra, let \(P\) be Lebesgue measure on \(\mathcal F_0\),
      and let \(\mathcal G=\{\emptyset,\Omega\}\). For every integrable
      \(X\), the constant \(EX\) is a version of \(E(X\mid\mathcal G)\),
      and every version equals \(EX\) at every point of \(\Omega\).
      The function \(Z'=EX+\mathbf1_{\{0\}}\) is \(\mathcal F_0\)-measurable
      and equals that version outside the \(\mathcal F_0\)-null set
      \(\{0\}\), but it is not \(\mathcal G\)-measurable and so is not a
      version.
\end{enumerate}
\end{hbtheorem}

Part~(iii) is the witness for the last sentence of
\cref{thm:p15-conditional-limits}: there a nonempty \(\mathcal F_0\)-null
set exists, and \(\mathcal G\) contains no nonempty null set at all.
Part~(i) is the kind of change used throughout the annex proofs of the
conditional limit theorems, where versions are reset on a null set that
belongs to \(\mathcal G\). Part~3 of \cref{thm:p15-conditional-examples}
is an instance of part~(i): there the value \(h(0)\) is changed on the
event \(\{Y=0\}\), which belongs to \(\sigma(Y)\).

\begin{hbtheorem}{thm:p15-w02b-p15-008-decreasing}{proved=annex:p15-w02b-p15-008-decreasing}{Conditional limits along decreasing sequences}
\begin{enumerate}
\item Let \(Y_1,Y_2,\ldots\) and \(Y\) be real
      \(\mathcal F_0\)-measurable variables with \(E|Y_1|<\infty\) and
      \(E|Y|<\infty\), and suppose that \(Y_n\downarrow Y\) almost
      everywhere: outside a null set, \(Y_n(\omega)\) is nonincreasing in
      \(n\) and converges to \(Y(\omega)\). Then every \(Y_n\) is
      integrable. If \(Z_n\) is a version of \(E(Y_n\mid\mathcal G)\) for
      each \(n\) and \(Z\) is a version of \(E(Y\mid\mathcal G)\), there is
      a set \(D\in\mathcal G\) with \(P(D)=0\) such that, for every
      \(\omega\notin D\), \(Z_n(\omega)\) is nonincreasing in \(n\) and
      converges to \(Z(\omega)\). In the notation of the chapter,
      \[
           C_{\mathcal G}Y_n\downarrow C_{\mathcal G}Y
           \qquad\text{almost everywhere.}
      \]
\item The hypothesis \(E|Y_1|<\infty\) cannot be removed. In the model of
      part~(iii) of \cref{thm:p15-w02b-p15-008-null-modification}, put
      \(Y_n(\omega)=\omega^{-1}\) for \(0<\omega<1/n\) and
      \(Y_n(\omega)=0\) otherwise, and \(Y=0\). Then
      \(Y_n\geq0\), \(Y_n\downarrow Y\) at every point and \(E|Y|=0\), but
      \(EY_n=\infty\) for every \(n\). The nonnegative extended version
      \(C_{\mathcal G}Y_n\) of \cref{thm:p15-conditional-limits} equals
      \(\infty\) at every point, while \(C_{\mathcal G}Y=0\).
\end{enumerate}
\end{hbtheorem}

Part~(i) is the remark that follows Durrett's Theorem~4.1.9, obtained
there by applying the increasing case to \(Y_1-Y_n\)
\cite[p.~210]{DurrettPTEAuthor2019}; the annex follows that route,
using item~1 of \cref{thm:p15-conditional-limits} as the increasing case.
The hypothesis \(E|Y|<\infty\) is what makes \(E(Y\mid\mathcal G)\)
defined in the sense of \cref{def:p15-general-conditional-expectation}.

\begin{hbexample}{def:p15-w02b-p15-008-decreasing-worked}{A decreasing sequence on a two-cell partition}
On \(\Omega=[0,1]\) with Borel sets and Lebesgue measure, let
\(\mathcal G\) be generated by the partition \(A_1=[0,\tfrac12]\),
\(A_2=(\tfrac12,1]\), and let \(Y_n(\omega)=\omega^n\). Then
\(Y_n\) decreases to \(Y=\mathbf1_{\{1\}}\) at every point, \(0\leq Y_n\leq1\),
and \(E|Y|=0\). By \cref{thm:p15-conditional-partitions}, a version of
\(E(Y_n\mid\mathcal G)\) is constant on each cell with value
\[
 2\int_0^{1/2}\omega^n\,d\omega=\frac{2^{-n}}{n+1}
 \quad\text{on }A_1,\qquad
 2\int_{1/2}^{1}\omega^n\,d\omega=\frac{2-2^{-n}}{n+1}
 \quad\text{on }A_2,
\]
and the constant \(0\) is a version of \(E(Y\mid\mathcal G)\). Both cell
values decrease in \(n\): the first plainly, and the second because
\((2-2^{-n-1})(n+1)\leq(2-2^{-n})(n+2)\) is equivalent to
\(2^{-n-1}(n+3)\leq2\), which holds for \(n\geq1\). Both tend to \(0\),
as part~(i) of \cref{thm:p15-w02b-p15-008-decreasing} requires.
\end{hbexample}

\section{Relationships}
\label{sec:p15-w02b-p15-008-relationships}

\Cref{def:p15-regular-conditional-law} is the special case of
\cref{def:p15-w02b-p15-008-rcd} with empty exceptional set, and
\cref{thm:p15-kernel-expectation} is then recovered from
\cref{thm:p15-w02b-p15-008-ae-integration}.
\Cref{thm:p15-w02b-p15-008-ae-reduction} is the bridge in the other
direction under a measurability hypothesis on the exceptional set, and
\cref{thm:p15-w02b-p15-008-ae-counterexample} shows that the hypothesis
cannot be dropped. The phenomenon behind the counterexample is the one
seen in part~3 of \cref{thm:p15-conditional-examples}: conditional
objects are determined only up to null sets, and a choice of values on a
null set can be compatible with every defining identity while failing a
further pointwise property. The joint-density kernel of
\cref{thm:p15-density-conditional-kernel} corresponds to the source's
continuation of its Example~4.1.6 \cite[p.~214]{DurrettPTEAuthor2019},
with the measure condition imposed at every argument.
